\documentclass[12pt,a4paper]{article}
\usepackage[utf8]{inputenc}
\usepackage[vietnamese]{babel}
\usepackage{amsmath,amssymb,amsthm}
\usepackage{geometry}
\usepackage{enumitem}
\usepackage[most]{tcolorbox}
\usepackage{hyperref}
\usepackage{fancyhdr}
\usepackage{tikz}
\usepackage{eso-pic}
\usepackage{xcolor}
\geometry{a4paper, margin=2cm, bottom=2.5cm}
\hypersetup{colorlinks=true, linkcolor=blue!70!black}
\setlist[itemize]{topsep=4pt, itemsep=2pt}
\setlist[enumerate]{topsep=4pt, itemsep=2pt}
\AddToShipoutPictureFG{%
  \AtPageCenter{%
    \begin{tikzpicture}[remember picture, overlay]
      \node[rotate=45, scale=1, opacity=0.06]{\fontsize{60}{72}\selectfont\bfseries\sffamily Đặng Tấn Quí};
    \end{tikzpicture}%
  }%
}
\pagestyle{fancy}
\fancyhf{}
\fancyhead[L]{\small\textit{HỌC MÁY}}
\fancyhead[R]{\small\textit{Đặng Tấn Quí}}
\fancyfoot[C]{\thepage}
\fancyfoot[L]{\footnotesize\textcopyright\ Đặng Tấn Quí}
\fancyfoot[R]{\footnotesize SĐT: 0377\,122\,210}
\renewcommand{\headrulewidth}{0.4pt}
\renewcommand{\footrulewidth}{0.4pt}
\fancypagestyle{plain}{\fancyhf{}\fancyfoot[C]{\thepage}\fancyfoot[L]{\footnotesize\textcopyright\ Đặng Tấn Quí}\fancyfoot[R]{\footnotesize SĐT: 0377\,122\,210}\renewcommand{\headrulewidth}{0pt}\renewcommand{\footrulewidth}{0.4pt}}
\newtcolorbox{dinhnghia}[1][]{enhanced, breakable, colback=blue!3!white, colframe=blue!60!black, fonttitle=\bfseries, title={Định nghĩa}, #1}
\newtcolorbox{dinhly}[1][]{enhanced, breakable, colback=green!3!white, colframe=green!50!black, fonttitle=\bfseries, title={Định lý}, #1}
\newtcolorbox{tinhchat}[1][]{enhanced, breakable, colback=yellow!5!white, colframe=orange!50!black, fonttitle=\bfseries, title={Tính chất}, #1}
\newtcolorbox{phuongphap}[1][]{enhanced, breakable, colback=gray!5!white, colframe=gray!50!black, fonttitle=\bfseries, title={Phương pháp}, #1}
\newtcolorbox{luuy}[1][]{enhanced, breakable, colback=red!3!white, colframe=red!50!black, fonttitle=\bfseries, title={Lưu ý}, #1}
\newtcolorbox{congthuc}[1][]{enhanced, breakable, colback=cyan!3!white, colframe=cyan!50!black, fonttitle=\bfseries, title={Công thức}, #1}
\newtcolorbox{vidu}[1][]{enhanced, breakable, colback=violet!3!white, colframe=violet!50!black, fonttitle=\bfseries, title={Ví dụ}, #1}

\begin{document}
\thispagestyle{empty}
\begin{tikzpicture}[remember picture, overlay]
  \fill[blue!80!black] (current page.south west) rectangle (current page.north east);
  \node[anchor=west, white, font=\fontsize{26}{32}\bfseries\selectfont]
    at ([xshift=3cm, yshift=-7.5cm]current page.north west) {HỌC MÁY};
  \node[anchor=west, white, font=\fontsize{18}{24}\selectfont]
    at ([xshift=3cm, yshift=-9cm]current page.north west) {Hồi quy, SVM, cây, neural network, K-Means, PCA};
  \node[anchor=west, white, font=\Large\bfseries] at ([xshift=3cm, yshift=-13cm]current page.north west) {Biên soạn:};
  \node[anchor=west, yellow!80!white, font=\fontsize{22}{28}\bfseries\selectfont] at ([xshift=3cm, yshift=-14.5cm]current page.north west) {ĐẶNG TẤN QUÍ};
  \node[anchor=south east, white, font=\Large\bfseries] at ([xshift=-2cm, yshift=3cm]current page.south east) {\the\year};
\end{tikzpicture}
\newpage
\tableofcontents
\newpage
%% ================================================================
\section{Nền tảng}

\begin{dinhnghia}[title={Phân loại bài toán}]
  \begin{itemize}
    \item \textbf{Học có giám sát}: dữ liệu $(x_i, y_i)$; hồi quy ($y \in \mathbb{R}$), phân loại ($y \in \{0,1,\ldots\}$).
    \item \textbf{Học không giám sát}: chỉ có $x_i$; phân cụm, giảm chiều.
    \item \textbf{Học tăng cường}: agent, môi trường, phần thưởng.
  \end{itemize}
\end{dinhnghia}

\begin{dinhnghia}[title={Bias--Variance tradeoff}]
  $\text{MSE} = \text{Bias}^2 + \text{Variance} + \text{Noise}$. Mô hình đơn giản: bias cao, variance thấp. Phức tạp: ngược lại.
\end{dinhnghia}

\begin{phuongphap}[title={Cross-validation}]
  $k$-fold CV: chia dữ liệu $k$ phần, luân phiên dùng 1 phần test, $k-1$ phần train. Ước lượng lỗi ngoài mẫu.
\end{phuongphap}

%% ================================================================
\section{Hồi quy}

\begin{dinhnghia}[title={Hồi quy tuyến tính}]
  $\hat{y} = \mathbf{w}^T\mathbf{x} + b$. Loss: $L = \frac{1}{n}\sum(y_i - \hat{y}_i)^2$. OLS: $\hat{\mathbf{w}} = (\mathbf{X}^T\mathbf{X})^{-1}\mathbf{X}^T\mathbf{y}$.
\end{dinhnghia}

\begin{dinhnghia}[title={Regularization}]
  Ridge ($L^2$): $\min \|\mathbf{y} - \mathbf{X}\mathbf{w}\|^2 + \lambda\|\mathbf{w}\|^2$.

  Lasso ($L^1$): $\min \|\mathbf{y} - \mathbf{X}\mathbf{w}\|^2 + \lambda\|\mathbf{w}\|_1$. Tạo sparsity (chọn biến).

  Elastic Net: kết hợp $L^1 + L^2$.
\end{dinhnghia}

%% ================================================================
\section{Phân loại}

\begin{dinhnghia}[title={Hồi quy logistic}]
  $P(y=1|\mathbf{x}) = \sigma(\mathbf{w}^T\mathbf{x} + b)$, $\sigma(z) = 1/(1+e^{-z})$. Loss: cross-entropy.
\end{dinhnghia}

\begin{dinhnghia}[title={k-NN (k Nearest Neighbors)}]
  Gán nhãn dựa trên đa số $k$ láng giềng gần nhất. Phi tham số, nhạy cảm với $k$ và metric.
\end{dinhnghia}

\begin{dinhnghia}[title={SVM (Support Vector Machine)}]
  $\max_{\mathbf{w}, b} \frac{2}{\|\mathbf{w}\|}$ t.c. $y_i(\mathbf{w}^T\mathbf{x}_i + b) \ge 1$.

  Kernel trick: $K(\mathbf{x}, \mathbf{x}') = \langle\phi(\mathbf{x}), \phi(\mathbf{x}')\rangle$. Kernel phổ biến: RBF $K = \exp(-\gamma\|\mathbf{x} - \mathbf{x}'\|^2)$.
\end{dinhnghia}

\begin{dinhnghia}[title={Cây quyết định}]
  Chia không gian đặc trưng bằng ngưỡng. Tiêu chí: Gini impurity $G = 1 - \sum p_k^2$, entropy $H = -\sum p_k\ln p_k$.

  \textbf{Random Forest}: tập hợp nhiều cây (bagging + random subspace).

  \textbf{Gradient Boosting}: xây cây tuần tự, mỗi cây fit residual. XGBoost, LightGBM.
\end{dinhnghia}

\begin{phuongphap}[title={Đánh giá phân loại}]
  Accuracy, Precision, Recall, F1-score. Ma trận nhầm lẫn. ROC curve, AUC.
\end{phuongphap}

%% ================================================================
\section{Mạng nơ-ron}

\begin{dinhnghia}[title={Perceptron nhiều lớp (MLP)}]
  $\mathbf{a}^{(l)} = \sigma(\mathbf{W}^{(l)}\mathbf{a}^{(l-1)} + \mathbf{b}^{(l)})$. Hàm kích hoạt: ReLU, sigmoid, tanh, softmax.
\end{dinhnghia}

\begin{phuongphap}[title={Backpropagation}]
  Chain rule: $\frac{\partial L}{\partial W^{(l)}} = \frac{\partial L}{\partial \mathbf{a}^{(l)}}\cdot\frac{\partial \mathbf{a}^{(l)}}{\partial W^{(l)}}$. Cập nhật: SGD, Adam.
\end{phuongphap}

\begin{tinhchat}[title={Kiến trúc phổ biến}]
  \begin{itemize}
    \item \textbf{CNN}: convolution, pooling --- hình ảnh.
    \item \textbf{RNN / LSTM / GRU}: chuỗi thời gian, NLP.
    \item \textbf{Transformer}: attention mechanism --- NLP, vision.
  \end{itemize}
\end{tinhchat}

%% ================================================================
\section{Học không giám sát}

\begin{phuongphap}[title={K-Means}]
  Phân $n$ điểm thành $k$ cụm. Lặp: gán cụm gần nhất, cập nhật centroid. Hội tụ (objective giảm).
\end{phuongphap}

\begin{phuongphap}[title={PCA}]
  Chiếu lên các thành phần chính (trị riêng lớn nhất của $\mathbf{X}^T\mathbf{X}$). Giảm chiều, trực quan hóa.
\end{phuongphap}

\begin{dinhnghia}[title={t-SNE, UMAP}]
  Phương pháp giảm chiều phi tuyến, bảo toàn cấu trúc cục bộ. Dùng để trực quan hóa dữ liệu chiều cao.
\end{dinhnghia}

\end{document}